\section{Methodology}
\begin{paracol}
    {2}

    This study follows factorial experimental design, analyzing the interplay between perception, coverage and system-level decision making.
    Two independent variables are defined:
    \begin{itemize}
    \item Detection method:
    
    \begin{itemize}
    \item Point cloud clustering
    \item 2D Depth projection
    \end{itemize}

    \item Coverage strategy:

    \begin{itemize}
    \item Boustrophedon Coverage
    \item Morphology based skeleton Coverage
    \item Spanning tree coverage
    \end{itemize}
    \end{itemize}

    This results in a total of 16 experimental configurations, however some of
    these setups are fundamentally incompatible, such as Systematic coverage with a
    clean as you go approach. The full setup is shown in table x.
    
    \medskip
    \centerline{\rule{7.8cm}{0.4pt}}
    \subsection{Evaluation Metrics}
    \subsubsection{Task Level Metrics}

    \begin{quote}
        Percentage of electronics, total completion time, failed attempts, distance traveled
    \end{quote}

    \subsubsection{Component Level Metrics}

    \begin{quote}
        Detection precision and recall, localization error, coverage percentage.
    \end{quote}

    \syncallcounters
    \switchcolumn

    .
    
\end{paracol}